\documentclass[10pt]{article}
\usepackage[T1]{fontenc}
\usepackage[utf8]{inputenc}
\usepackage[french]{babel}
\usepackage{amsmath,amssymb,mathrsfs,mathpazo}
\usepackage{geometry}
\geometry{a4paper,margin=2cm}
\usepackage{graphicx,array,tabularx,tabularray,colortbl,multirow,multicol,booktabs,enumitem,pifont,ulem,fancybox,diagbox,cancel}
\usepackage{pstricks,pstricks-add,pst-plot,pst-tree,pst-node,pst-text,pst-func,pst-math,pst-bspline,pst-3dplot,pst-eucl,pst-solides3d}
\let\pscalculate\undefined
\usepackage{pst-bezier}
\usepackage{tikz,pgfplots,tkz-tab}
\usetikzlibrary{arrows,arrows.meta,patterns,calc,positioning,trees,decorations.markings,shapes.geometric,shapes.misc,shapes.arrows,intersections,decorations.pathmorphing,decorations.pathreplacing,backgrounds,shadows,babel}
\pgfplotsset{compat=1.18}
\usepgfplotslibrary{fillbetween}
\usepackage[np]{numprint}
\usepackage{listings,makecell,pas-tableur,twoopt}
\let\barre\undefined
\usepackage[pstricks]{tablvar}
\definecolor{darkgreen}{RGB}{0,100,0}
\providecommand{\red}{\color{red}}
\providecommand{\blue}{\color{blue}}
\providecommand{\green}{\color{green}}
\providecommand{\black}{\color{black}}
\providecommand{\white}{\color{white}}
\providecommand{\gray}{\color{gray}}
\providecommand{\cyan}{\color{cyan}}
\providecommand{\magenta}{\color{magenta}}
\providecommand{\yellow}{\color{yellow}}
\providecommand{\violet}{\color{violet}}
\providecommand{\orange}{\color{orange}}
\providecommand{\euro}{\texteuro}
\providecommand{\EUR}{\texteuro}
\providecommand{\ieme}{\textsuperscript{e}}
\providecommand{\degres}{\ensuremath{^\circ}}
\providecommand{\degree}{\ensuremath{^\circ}}
\providecommand{\vv}[1]{\overrightarrow{#1}}
\providecommand{\Coord}[1]{#1}
\providecommand{\decosix}{}
\providecommand{\starredbullet}{\ensuremath{\bullet}}
\pagestyle{empty}
\setlength{\parindent}{0pt}
\usepackage{adjustbox}\DeclareUnicodeCharacter{2236}{:}\DeclareUnicodeCharacter{2190}{\ensuremath{\leftarrow}}
\providecommand{\R}{\mathbb{R}}
\providecommand{\N}{\mathbb{N}}
\providecommand{\Z}{\mathbb{Z}}
\providecommand{\Q}{\mathbb{Q}}
\providecommand{\D}{\mathbb{D}}
\providecommand{\C}{\mathbb{C}}
\providecommand{\e}{{\mathrm{e}}}
\providecommand{\vect}[1]{\overrightarrow{#1}}
\providecommand{\vectt}[1]{\overrightarrow{\mathrm{#1}}}
\providecommand{\vvt}[1]{\overrightarrow{\mathrm{#1}}}
\providecommand{\np}[2][]{{#2~#1}}
\providecommand{\vv}[1]{{\overrightarrow{#1}}}
\renewcommand{\d}{{\mathrm{d}}}
\providecommand{\degres}{{^\circ}}
\providecommand{\degre}{{^\circ}}
\providecommand{\blue}{\color{blue}}
\providecommand{\red}{\color{red}}
\providecommand{\enskip}{\quad}
\providecommand{\llbracket}{\lbrack\!\lbrack}
\providecommand{\rrbracket}{\rbrack\!\rbrack}
\providecommand{\euro}{{\text{€}}}
\providecommand{\no}{{\text{n°}}}
\providecommand{\interval}[2]{{[#1;#2]}}
\providecommand{\intervFO}[2]{{[#1;#2[}}
\providecommand{\dotfill}{\ldots}

\providecommand{\ds}{\displaystyle}
\providecommand{\textding}[1]{#1}
\providecommand{\Oij}{\(\mathrm{O};\vec{\imath},\vec{\jmath}\)}
\providecommand{\Oijk}{\(\mathrm{O};\vec{\imath},\vec{\jmath},\vec{k}\)}
\providecommand{\Ouv}{\(\mathrm{O};\vec{u},\vec{v}\)}
\providecommand{\intervFF}[2]{\ensuremath{[#1;#2]}}
\renewcommand{\euro}{\ifmmode\text{\texteuro}\else\texteuro\fi}
\DeclareUnicodeCharacter{00B0}{\ifmmode^{\circ}\else\textdegree\fi}
\DeclareUnicodeCharacter{20AC}{\ifmmode\text{\texteuro}\else\texteuro\fi}
% Énoncés : Éducation nationale. Transcription et corrigés : APMEP, Denis Vergès et contributeurs. Découpage et classement : CoopMaths, Sébastien Lozano, Éric Elter et contributeurs. Adaptation de la présentation : Maths Explorer.
% CC BY-SA : https://coopmaths.fr/www/about/
% https://coopmaths.fr/alea/static/bac/2025/tex/bac_2025_06_sujet1_polynesie_2.tex
\begin{document}
\noindent\textit{APMEP / CoopMaths. CC BY-SA. Présentation adaptée par Maths Explorer.}\par\medskip
\medskip

Deux avions sont en approche d'un aéroport.

On munit l'espace d'un repère orthonormé \Oijk dont J'origine O est le pied de la
tour de contrôle, et le sol est le plan $P_0$ d'équation $z = 0$.

L'unité des axes correspond à 1 km.

On modélise les avions par des points.

\begin{center}
\psset{unit=1cm,arrowsize=2pt 3}
\begin{pspicture}(-6,-1.8)(4.7,5.5)

\psline{->}(-6,0.3)(4.7,-0.1)
\psline{->}(0,-1.8)(0,5.5)
\psline{->}(-3,-1)(3,1)
\multido{\n=-1.3+1.3,\nb=-2+2}{5}{\rput(0,\n){\small $\bullet$}\uput[l](0,\n){\small \nb}}
\multido{\n=-4.8+1.2,\na=0.25 +-0.05,\nb=-8+2}{8}{\rput(\n,\na){\small $\bullet$}\uput[d](\n,\na){\small \nb}}
\multido{\n=-2.4+0.6,\na=-0.8 +0.2,\nb=-8+2}{9}{\rput(\n,\na){\small $\bullet$}\uput[ul](\n,\na){\small \nb}}
\psline[linewidth=1.5pt](-4.1,5.5)(-1.7,-0.3)\psline[linewidth=1.5pt,linestyle=dashed](-1.7,-0.3)(-1.4,-1)
\psline[linewidth=1.5pt](-6,5.5)(1,-0.4)\psline[linewidth=1.5pt,linestyle=dashed](1,-0.4)(1.8,-1)
\uput[d](4.6,-0.15){$x$}\uput[d](0,5.3){$z$}\uput[ul](3,1){$y$}
\uput[ur](-3.5,4){$d_{\text{A}}$}\uput[dl](-4.2,4){$d_{\text{B}}$}
\rput(-3.9,5){$\bullet$}\rput(-4.8,4.5){$\bullet$}
\uput[ur](-3.9,4.95){A}\uput[dl](-4.8,4.5){B}
\end{pspicture}
\end{center}

L'avion Alpha transmet à la tour sa position en A$(-7~;~ 1~;~7)$ et sa trajectoire est dirigée par le vecteur $\vect{u}\begin{pmatrix}2\\-1\\-3\end{pmatrix}$.

L'avion Bêta transmet une trajectoire définie par la droite $d_{\text{B}}$ passant par le point B dont une représentation paramétrique est :

\[\left\{\begin{array}{l c l}
x &=&-11 + 5t\\
y &=&-5 + \phantom{5}t\\
z &=&11 - 4t \end{array}\right.\:\text{où}\: t\:\text{décrit}\:\: \R\]

\medskip

\begin{enumerate}
\item S'il ne dévie pas de sa trajectoire, déterminer les coordonnées du point S en lequel l'avion Bêta touchera le sol.
\item 
	\begin{enumerate}
		\item Déterminer une représentation paramétrique de la droite $d_{\text{A}}$ caractérisant la trajectoire de l'avion Alpha.
		\item Les deux avions peuvent-ils entrer en collision ?
	\end{enumerate}
\item 
	\begin{enumerate}
		\item Démontrer que l'avion Alpha passe par la position E$(-3~;~-1~;~1)$.
		\item Justifier qu'une équation cartésienne du plan $P_{\text{E}}$ passant par E et perpendiculaire à la droite $d_{\text{A}}$ est:

\[2x - y - 3z + 8 = 0.\]
		\item Vérifier que le point F$(-1~;~-3~;~3)$ est le point d'intersection du plan $P_{\text{E}}$ et de la droite $d_{\text{B}}$.
		\item Calculer la valeur exacte de la distance EF, puis vérifier que cela correspond à une distance de \np{3464} m, à 1 m près.
	\end{enumerate}
\item La réglementation aérienne stipule que deux avions en approche doivent être à tout instant à au moins 3 milles nautiques l'un de l'autre (1 mille nautique vaut \np{1852}~m).

Si les avions Alpha et Bêta sont respectivement en E et F au même instant, leur distance de sécurité est-elle respectée ?
\end{enumerate}

\bigskip

\end{document}
